% ============================================================
% MODULO 05 - METODOLOGIA (redacao revisada; NBR 6022:2018)
% ============================================================
\section{Metodologia}

\subsection{Desenho}

Adotou-se uma revis\~ao integrativa de car\'ater narrativo, acompanhada de vinheta cl\'inica \'unica e anonimizada. O trabalho n\~ao \'e uma revis\~ao sistem\'atica: n\~ao segue o protocolo PRISMA em sua integralidade e n\~ao realiza meta-an\'alise. A vinheta cumpre fun\c{c}\~ao ilustrativa e heur\'istica, sem qualquer pretens\~ao inferencial. Cabe justificar essa escolha: um caso \'unico serve para gerar hip\'oteses, n\~ao para confirm\'a-las. Por isso o desenho dispensa, deliberadamente, grupo de controle, amostragem probabil\'istica e isolamento de vari\'aveis -- exig\^encias que s\'o fariam sentido se o objetivo fosse testar efic\'acia, o que n\~ao \'e o caso.

\subsection{Crit\'erios de invalida\c{c}\~ao da hip\'otese}

Para que a hip\'otese do ``tempo verde-social'' seja pass\'ivel de teste, registram-se as condi\c{c}\~oes que a enfraqueceriam ou refutariam: (a) SNAP-IV com piora ou estabilidade dos sintomas ap\'os tr\^es meses de brincar regular na vizinhan\c{c}a; (b) manuten\c{c}\~ao da melhora mesmo com interrup\c{c}\~ao prolongada do brincar social, o que apontaria matura\c{c}\~ao; (c) surgimento de registro de psicof\'armaco iniciado, ajustado ou suspenso no per\'iodo; (d) discord\^ancia entre informantes (m\~ae e professora) na escala padronizada; (e) melhora equivalente em crian\c{c}a do mesmo contexto sem exposi\c{c}\~ao ao brincar de vizinhan\c{c}a.

\subsection{Fontes e estrat\'egia de busca}

As buscas explorat\'orias ocorreram em outubro de 2026, nas bases SciELO, BVS/PEPSIC, Portal de Peri\'odicos CAPES, PubMed e PsycINFO, cobrindo o per\'iodo de 2004 a 2026, nos idiomas portugu\^es, ingl\^es e espanhol. Em portugu\^es (DeCS), combinaram-se: ``Transtorno do D\'eficit de Aten\c{c}\~ao com Hiperatividade'' AND ``Pr\'e-escolar'' AND ``Jogos e Brinquedos'' AND ``Meio Ambiente'' AND ``Terapia Centrada no Cliente''. Em ingl\^es (MeSH): ``Attention Deficit Hyperactivity Disorder'' AND ``Preschool Child'' AND ``Play'' AND ``Green Space/Neighborhood'' AND ``Person-Centered Therapy/Play Therapy''. A string final combinou os termos ``ADHD'' ou ``TDAH''; ``preschool'' ou ``pr\'e-escolar''; ``free play'', ``outdoor play'', ``neighborhood'', ``brincar livre'' ou ``vizinhan\c{c}a''; e ``nature'', ``green space'' ou ``natureza''. Entraram estudos com crian\c{c}as de 3 a 12 anos, com laudo ou tra\c{c}os de TDAH, em contexto externo, de vizinhan\c{c}a ou de natureza, e desfecho comportamental de aten\c{c}\~ao, hiperatividade ou intera\c{c}\~ao social. Foram exclu\'idos trabalhos s\'o com adultos, farmacologia sem componente ambiental, resumos sem texto completo e duplicatas. Os dados foram organizados em matriz com refer\^encia, desenho, idade, contexto, medida, achado e limite.

\subsection{Cuidados \'eticos da vinheta}

A crian\c{c}a recebeu nome fict\'icio (M.); endere\c{c}o, escola, nomes de familiares, imagens e datas identificadoras foram suprimidos. O relato exige Termo de Consentimento Livre e Esclarecido do respons\'avel e assentimento l\'udico da crian\c{c}a [A ANEXAR -- se a pesquisa for institucional, informar o n\'umero do Comit\^e de \'Etica em Pesquisa; se for relato de experi\^encia cl\'inica para fins did\'aticos, declarar o consentimento para uso anonimizado, conforme o C\'odigo de \'Etica Profissional do Psic\'ologo, o Estatuto da Crian\c{c}a e do Adolescente (Lei n. 8.069/1990) e a Lei Geral de Prote\c{c}\~ao de Dados (Lei n. 13.709/2018)]. As fontes do caso foram anamnese materna, prontu\'ario, observa\c{c}\~ao em [N] sess\~oes de 50 minutos [A PREENCHER] e di\'ario de campo. Nenhuma escala padronizada foi aplicada at\'e o fechamento -- limita\c{c}\~ao declarada --, recomendando-se SNAP-IV (pais e professores) e CBCL 1,5--5 na continuidade do acompanhamento.

\subsection{Medicamento n\~ao declarado}

At\'e o fechamento deste manuscrito, a fam\'ilia n\~ao declarou uso atual, pregresso, dose, ades\~ao ou suspens\~ao de psicof\'armacos, nem outros tratamentos concomitantes (psicopedagogia, terapia ocupacional, interven\c{c}\~oes escolares). O dado foi registrado como ``informa\c{c}\~ao n\~ao declarada'', sem que se presuma aus\^encia de medica\c{c}\~ao. Essa lacuna impede qualquer infer\^encia farmacol\'ogica e \'e retomada entre as limita\c{c}\~oes.

\subsection{Confundidores n\~ao controlados}

Foram tratadas como explica\c{c}\~oes alternativas: matura\c{c}\~ao neurol\'ogica entre 5,0 e 5,5 anos; mudan\c{c}as de sono, rotina e alimenta\c{c}\~ao; nova escola ou educadora; redu\c{c}\~ao do estresse familiar ap\'os a mudan\c{c}a; sazonalidade; tempo de tela antes e depois; amadurecimento da pr\'opria alian\c{c}a terap\^eutica; e vi\'es de expectativa materna.
